\begin{exercise}{80}
Show that a functor $\mathcal C\to\mathcal D$ between $\infty$-categories is an
isofibration if and only if $\mathcal C^{\mathrm{op}}\to\mathcal D^{\mathrm{op}}$
is an isofibration.

\begin{proof}
The ordinary isofibration condition in the proof of
Exercise~\ref{ex:79} is invariant under taking opposites.
Indeed, let $q\colon A\to B$ satisfy that condition.
For an isomorphism $\beta\colon b\to q(a)$, lift
$\beta^{-1}\colon q(a)\to b$ to an isomorphism
$\alpha\colon a\to a'$ with $q(a')=b$ and
$q(\alpha)=\beta^{-1}$.
Then $\alpha^{-1}\colon a'\to a$ lies over $\beta$.
These incoming isomorphism lifts are precisely the outgoing
isomorphism lifts required for $q^{\mathrm{op}}$.
The converse follows by taking opposites again.

Definition~1.2.5 gives a natural isomorphism
$h(\mathcal C^{\mathrm{op}})\cong(h\mathcal C)^{\mathrm{op}}$:
it is the identity on objects and sends the class of a reversed
edge to the opposite of its class.
Reversing edges reverses the triangle composition relations
and preserves the identity relations.
The additional homotopy relations already follow from
triangles with an identity edge.
Thus the assignment is well-defined and inverse to the
same construction.
It is natural in functors because both assignments apply
the functor to each edge generator.

Now let $p\colon\mathcal C\to\mathcal D$ be an isofibration.
By Definition~2.1.4 it is an inner fibration.
Exercises~\ref{ex:48} and~\ref{ex:51}(2) show that
$p^{\mathrm{op}}$ is an inner fibration between
$\infty$-categories.
Exercise~\ref{ex:79} says that $hp$ satisfies the ordinary
isofibration condition.
By the preceding paragraphs, so does $h(p^{\mathrm{op}})$,
identified naturally with $(hp)^{\mathrm{op}}$.
Applying Exercise~\ref{ex:79} to $p^{\mathrm{op}}$ proves
that it is an isofibration.
Applying this implication to $p^{\mathrm{op}}$ proves the converse.
\end{proof}
\end{exercise}
